\documentclass[11pt]{article}
\usepackage[margin=1in]{geometry}
\usepackage{setspace}
\usepackage{hyperref}
\usepackage{natbib}
\bibliographystyle{chicago}

\title{Component 2: Summary of Research Findings}
\author{Alejandro De La Torre}
\date{March 2026}

\begin{document}
\maketitle

\onehalfspacing

This paper examines an important financial risk facing older households: the income shock that can occur when one spouse dies. In the United States, many married couples depend heavily on Social Security benefits in retirement. When a spouse dies, however, the household typically loses one of those benefit checks, even though many of the household’s expenses remain unchanged. As a result, surviving spouses—most often widows—can experience a sharp decline in financial resources. The research paper by Slavov asks whether the timing of Social Security claiming decisions made earlier in retirement can reduce the financial strain that surviving spouses face after widowhood.\footnote{Sita N. Slavov, “The Impact of Spousal Social Security Claiming Decisions on the Financial Shock of Widowhood,” Working Paper No. 34612 (Cambridge, MA: National Bureau of Economic Research, 2025), https://doi.org/10.3386/w34612.}

The study focuses on how the claiming behavior of the primary earner—typically the husband in older cohorts—affects the financial well-being of the surviving spouse. Under Social Security rules, workers receive larger monthly payments if they delay claiming benefits beyond the earliest eligibility age. These higher benefits also raise the survivor benefit that a spouse can receive after the worker dies. Because widows are generally entitled to the larger of their own benefit or their deceased spouse’s benefit, delayed claiming by the husband can increase the income available to the widow later in life.

To study this issue, the author uses survey data that track older Americans’ Social Security claiming behavior, marital histories, and income. The analysis compares households in which husbands claim benefits early with those in which husbands delay claiming benefits until full retirement age or later. Rather than focusing only on average income levels in retirement, the study examines the change in household income that occurs when widowhood begins. This approach allows the author to measure more precisely the financial shock that surviving spouses experience.

The results show that delayed claiming significantly reduces the financial shock associated with widowhood. When husbands postpone claiming Social Security benefits, their higher monthly payments translate into larger survivor benefits for their wives. As a result, widows experience a smaller drop in income after their spouse dies. In effect, a decision made earlier in retirement—when to begin collecting Social Security—can shape financial security many years later.

These findings carry important implications for retirement planning and policy. Many households claim benefits as soon as they become eligible, often due to health concerns or immediate financial needs. However, early claiming can leave surviving spouses financially vulnerable. Encouraging households to consider survivor protection when making claiming decisions could therefore improve long-term financial security for widows.

More broadly, the paper highlights the role of Social Security not only as a retirement income program but also as a form of survivor insurance. Claiming decisions are often framed as individual choices, yet they can shape the financial stability of an entire household. By showing that delayed claiming can cushion the economic impact of widowhood, the study contributes to policy discussions about how Social Security rules influence retirement behavior and financial well-being later in life.

\bibliography{../brookings}
\end{document}